\section{Final assembly and the full root-separation constant}
\label{sec:final-assembly-near-root-v4}

We combine the chromatic lower location, the one-part-buffer signed profile,
and the quantitative amplification estimate. The final high-probability events
are intersected by a union bound; no independence is used.

\subsection{The phase-resolved full gap}

Let $r_+(n)$ be the ordinary first-moment root and let
$r_4^{\mathrm{co}}(n)$ be the signed four-size root. Section~5 gives
\begin{equation}
  r_+(n)-r_4^{\mathrm{co}}(n)
  =
  \left[
    \frac{(\log 2)^2}{4}A_4(\delta_n)+o(1)
  \right]
  \frac{n}{(\log n)^3},
  \qquad
  A_4(\delta):=\log 2-D_4(\delta).
  \label{eq:phase-resolved-root-gap-v4}
\end{equation}
The signed witness is placed at
\[
  k_{\mathrm{co}}
  =\left\lceil r_4^{\mathrm{co}}(n)\right\rceil+1.
\]
Thus $k_{\mathrm{co}}-r_4^{\mathrm{co}}\in[1,2)$, and the bounded placement
error is negligible on the scale $n/(\log n)^3$. Consequently
\begin{equation}
  r_+(n)-k_{\mathrm{co}}
  =
  \left[
    \frac{(\log 2)^2}{4}A_4(\delta_n)+o(1)
  \right]
  \frac{n}{(\log n)^3}.
  \label{eq:near-root-retained-gap-v4}
\end{equation}
Unlike midpoint placement, the one-part buffer retains the complete leading
root separation. Its only purpose is to make the signed first moment uniformly
larger than one. Equation (5.19) gives the stronger quantitative statement
\[
  \log\mathbb E Z_{\mathbf k}^{\mathrm{sgn}}
  =\Theta((\log n)^2)
\]
uniformly in the phase.

Section~4 constructs a deterministic integer $k_\chi^-$ such that
\[
  \mathbb P\bigl(\chi(G_n)>k_\chi^-\bigr)\longrightarrow1.
\]
Equation (5.20) gives
\begin{equation}
  k_\chi^- - k_{\mathrm{co}}
  =
  \left[
    \frac{(\log 2)^2}{4}A_4(\delta_n)+o(1)
  \right]
  \frac{n}{(\log n)^3}.
  \label{eq:integer-full-location-gap-v4}
\end{equation}

The sharpened Section~9 estimate supplies a phase-independent constant
$C_\sharp$ for which
\[
  \frac{\mathbb E Z^2}{(\mathbb E Z)^2}
  \le
  \exp\!\left(C_\sharp\sqrt n\,(\log n)^{3/2}\right).
\]
Put
\[
  \Lambda_n^\sharp
  :=C_\sharp\sqrt n\,(\log n)^{3/2}.
\]
Paley--Zygmund gives
\[
  \mathbb P(Z>0)\ge e^{-\Lambda_n^\sharp}.
\]
Apply Lemma~10.2 with $k_n=k_{\mathrm{co}}$ and
$r_n=\log n$.  The number of additional cocoloring classes is bounded by
\[
\begin{aligned}
  a_n
  &:=C\left(
    \frac{\sqrt{n\Lambda_n^\sharp}+\sqrt{nr_n}}{\log n}
    +n^{1/3}+1
  \right)\\
  &=O\!\left(n^{3/4}(\log n)^{-1/4}\right).
\end{aligned}
\]
Indeed,
\[
  \frac{\sqrt{n\Lambda_n^\sharp}}{\log n}
  =O\!\left(n^{3/4}(\log n)^{-1/4}\right),
  \qquad
  \frac{\sqrt{nr_n}}{\log n}
  =O\!\left(\sqrt{\frac n{\log n}}\right).
\]
Moreover,
\[
  \frac{n^{3/4}(\log n)^{-1/4}}{n/(\log n)^3}
  =\frac{(\log n)^{11/4}}{n^{1/4}}
  \longrightarrow0.
\]
Therefore
\begin{equation}
  a_n=o\!\left(\frac{n}{(\log n)^3}\right),
  \qquad
  \mathbb P\bigl(\zeta(G_n)>k_{\mathrm{co}}+a_n\bigr)
  \longrightarrow0.
  \label{eq:near-root-amplification-v4}
\end{equation}
This explicit amplification loss is much smaller than the root separation.

Intersecting the event in \eqref{eq:near-root-amplification-v4} with the
chromatic lower event and using
\eqref{eq:integer-full-location-gap-v4}, we obtain a deterministic
$\epsilon_n\to0$ such that
\begin{equation}
  \mathbb P\!\left(
    \chi(G_n)-\zeta(G_n)
    \ge
    \left[
      \frac{(\log 2)^2}{4}A_4(\delta_n)-\epsilon_n
    \right]
    \frac{n}{(\log n)^3}
  \right)
  \longrightarrow1.
  \label{eq:phase-resolved-full-final-v4}
\end{equation}

\subsection{The exact four-support certificate and fixed slack}

It remains to replace $A_4(\delta)$ by a phase-independent constant while
retaining fixed positive slack. Put $q=\log 2$, let
\[
  w_i(\lambda)=\exp\!\left(\lambda i-\frac q2i^2\right),
  \qquad
  Z_4(\lambda)=\sum_{i=2}^{5}w_i(\lambda),
\]
and write
\[
  M_4(\lambda)
  =
  \frac{\sum_{i=2}^{5}i w_i(\lambda)}{Z_4(\lambda)}.
\]
For the omitted low and high parts of the unrestricted support, define
\[
  L(\lambda)
  :=
  \frac{\sum_{i=-1}^{1}w_i(\lambda)}{Z_4(\lambda)},
  \qquad
  H(\lambda)
  :=
  \frac{\sum_{i=6}^{\infty}w_i(\lambda)}{Z_4(\lambda)}.
\]
The mean $M_4$ is strictly increasing, $L$ is strictly decreasing, and $H$ is
strictly increasing. The derivative of each logarithmic partition ratio is the
difference of the corresponding tilted means, and the supports are ordered as
\[
  \{-1,0,1\}<\{2,3,4,5\}<\{6,7,\ldots\}.
\]

The following inequalities are exact rational certificates:
\[
\begin{array}{c|c|c}
\lambda & \text{mean certificate} & \text{omitted-mass certificate}\\
\hline
\frac{49}{20}q
  & M_4(\lambda)<\frac{2}{q}
  & L(\lambda)<\frac{2629}{10000}\\[2mm]
\frac{29}{10}q
  & {}
  & L(\lambda)<\frac{329}{2500},\quad
    H(\lambda)<\frac{37}{2500}\\[2mm]
\frac{83}{20}q
  & M_4(\lambda)>1+\frac{2}{q}
  & H(\lambda)<\frac{357}{2500}.
\end{array}
\]
For completeness, set
\[
  r=2^{1/20},
  \qquad
  r_-:=\frac{1035264923841377}{10^{15}},
  \qquad
  r_+:=\frac{1035264923841378}{10^{15}}.
\]
The integer inequalities $r_-^{20}<2<r_+^{20}$ give
$r_-<r<r_+$. At each tilt in the table, every finite weight is an integral
power of $r$, so upper and lower bounds follow by substituting $r_+$ or $r_-$
according to the sign of that power. For the two high tails, sum the first term
and bound the remainder geometrically; the ratio $w_{i+1}/w_i$ decreases by a
factor $1/2$ when $i$ increases by one. Finally,
\[
  \sum_{k=1}^{N}\frac{1}{k2^k}
  <q<
  \sum_{k=1}^{N}\frac{1}{k2^k}
  +\frac{1}{(N+1)2^N}
\]
with $N=200$ supplies rational bounds for $q$. The companion script
\texttt{check\_constant\_ledger\_v3.py} performs these integer comparisons.

Let $\lambda_4$ be the tilt of the four-support optimizer at
$T_0=1+2/q-\delta$. Since $2/q\le T_0\le1+2/q$, the mean certificates imply
\[
  \frac{49}{20}q<\lambda_4<\frac{83}{20}q.
\]
If $\lambda_4\le(29/10)q$, monotonicity gives
\[
  L(\lambda_4)+H(\lambda_4)
  <
  \frac{2629}{10000}+\frac{37}{2500}
  =
  \frac{2777}{10000}.
\]
If $\lambda_4\ge(29/10)q$, then
\[
  L(\lambda_4)+H(\lambda_4)
  <
  \frac{329}{2500}+\frac{357}{2500}
  =
  \frac{2744}{10000}
  <
  \frac{2777}{10000}.
\]
Thus the same strict rational bound holds throughout the phase.

The unrestricted partition function at $\lambda_4$ is
\[
  Z_+(\lambda_4)
  =Z_4(\lambda_4)
   \bigl(1+L(\lambda_4)+H(\lambda_4)\bigr).
\]
Evaluating the unrestricted dual objective at the four-support optimizer gives
\[
  D_4(\delta)
  \le
  \log\bigl(1+L(\lambda_4)+H(\lambda_4)\bigr)
  <
  \log\!\left(\frac{12777}{10000}\right).
\]
Consequently, uniformly for $0\le\delta\le1$,
\begin{equation}
  A_4(\delta)
  >
  \log\!\left(\frac{20000}{12777}\right).
  \label{eq:exact-four-support-certificate-v4}
\end{equation}
The certificate has a fixed margin over the constant stated below:
\begin{equation}
  \log\!\left(\frac{20000}{12777}\right)
  =
  \log\!\left(\frac{1000}{639}\right)
  +
  \underbrace{\log\!\left(\frac{12780}{12777}\right)}_{=:~\sigma_4>0}.
  \label{eq:fixed-four-support-slack-v4}
\end{equation}

Combining \eqref{eq:phase-resolved-full-final-v4} with
\eqref{eq:exact-four-support-certificate-v4}, and then taking $n$ large enough
that
\[
  \epsilon_n<\frac{(\log 2)^2}{4}\sigma_4,
\]
yields
\[
  \mathbb P\!\left(
    \chi(G_n)-\zeta(G_n)
    \ge
    \frac{(\log 2)^2}{4}
    \log\!\left(\frac{1000}{639}\right)
    \frac{n}{(\log n)^3}
  \right)
  \longrightarrow1.
\]
The displayed coefficient equals
\[
  \frac{(\log 2)^2}{4}
  \log\!\left(\frac{1000}{639}\right)
  =0.053792819616758\ldots.
\]
Before absorbing the asymptotic error, the stronger certificate gives
\[
  \frac{(\log 2)^2}{4}
  \log\!\left(\frac{20000}{12777}\right)
  =0.053821018526027\ldots.
\]

\ifErdosProofClosed
This proves the main theorem. Since $n/(\log n)^3\to\infty$, the
chromatic--cochromatic difference tends to infinity with high probability
along the full sequence of integers.
\else
This proves the strengthened conditional verification statement, subject to
the closure gates listed in the appendices. The publication switch remains
disabled until the one-part-buffer interface and the preceding proof packages
have been independently replayed together.
\fi

\begin{corollary}[Simultaneous complement form]
Under the hypotheses used above,
\[
  \min\{\chi(G_n),\chi(\overline{G_n})\}-\zeta(G_n)
  \ge
  \frac{(\log 2)^2}{4}
  \log\!\left(\frac{1000}{639}\right)
  \frac{n}{(\log n)^3}
\]
with high probability.
\end{corollary}

\begin{proof}
The complement $\overline{G_n}$ again has law $G(n,1/2)$, and
$\zeta(\overline G)=\zeta(G)$. Apply the preceding result to $G_n$ and to
$\overline{G_n}$ and intersect the two events by a union bound. No further
asymptotic loss is introduced.
\end{proof}
